\documentclass[11pt]{article}
% Latin Modern with T1 encoding when available; otherwise Computer Modern (OT1).
\IfFileExists{lmodern.sty}{\usepackage{lmodern}\usepackage[T1]{fontenc}}{}
\usepackage[utf8]{inputenc}
\usepackage{microtype}
\usepackage{amsmath,amssymb,amsthm,mathtools,booktabs,array}
\usepackage[margin=1in]{geometry}
\usepackage{xcolor}
\usepackage{float}
\usepackage{listings}
\usepackage[hyphens]{url}
\usepackage{hyperref}
\hypersetup{colorlinks=true,linkcolor=blue!45!black,citecolor=blue!45!black,
  urlcolor=blue!45!black,
  pdftitle={Integer multiplication below n log n: the exponent kappa = 2\textasciicircum{}-74}}

\numberwithin{equation}{section}
\numberwithin{table}{section}
\newtheorem{theorem}{Theorem}[section]
\newtheorem{lemma}[theorem]{Lemma}
\newtheorem{proposition}[theorem]{Proposition}
\newtheorem{corollary}[theorem]{Corollary}
\theoremstyle{remark}
\newtheorem{remark}[theorem]{Remark}

\newcommand{\OAI}{\cite{OAI}}
\newcommand{\Wb}{W_{\mathrm b}}
\newcommand{\sbb}{s_{\mathrm b}}
\newcommand{\Lb}{L_{\mathrm b}}
\newcommand{\zb}{z_{\mathrm b}}
\newcommand{\etab}{\eta_{\mathrm b}}
\newcommand{\Wc}{W_{\mathrm c}}
\newcommand{\scc}{s_{\mathrm c}}
\newcommand{\lean}[1]{\texttt{#1}}

\lstset{basicstyle=\ttfamily\scriptsize,columns=fullflexible,keepspaces=true,
  breaklines=true,frame=single,framerule=0.3pt,xleftmargin=2pt,xrightmargin=2pt,
  showstringspaces=false,upquote=true}

\title{Integer multiplication below \texorpdfstring{$n\log n$}{n log n}:\\
the exponent \texorpdfstring{$\kappa=2^{-74}$}{kappa = 2\^{}-74}}
\author{}
\date{October 7, 2026}

\begin{document}
\maketitle

\begin{abstract}
The multitape algorithm of the preprint \emph{Integer multiplication
below $n\log n$} \OAI{} multiplies two $n$-bit integers in time
$O(n(\lg n)^{1-\kappa})$ with $\kappa=2^{-182}$; its constants were
chosen conservatively. We show that the same algorithm, with two local
changes and new constants, achieves $\kappa=2^{-74}$. First, the axes of
the Gaussian-resampling and Chinese-remainder layouts can be moved with
$O(d)$ direct chunk interchanges instead of $O(d^2)$ adjacent ones. This
removes the condition $\epsilon(2-\tau)<1-\tau$ that forced the
dimension exponent $\epsilon$ below $2^{-50}$. Second, the
address-interchange network of \OAI{} with ground set of size $46$
instead of $100$ admits the swap exponent $\tau=1-\tfrac65 2^{-36}$, against
the $1-2^{-50}$ used in \OAI. These two changes alone give
$\kappa=2^{-76}$. An additive unrolling of the guard-width recursion
and a smaller Gaussian width then allow $\epsilon=3/16$ and give
$\kappa=2^{-74}$. We also show that, with any bit network of the type
used in \OAI, every admissible choice of constants in this architecture
has $\kappa<2^{-73.7}$, so further progress needs a new idea rather than
tuning. All conditions and margins are checked in exact rational
arithmetic and, independently, in a companion Lean~4 file; the scan over
network sizes and the barrier constant are evaluated numerically.
\end{abstract}

\tableofcontents

\section{Introduction}

Throughout, \OAI{} denotes the 73-page preprint dated September 23,
2026, and references such as ``Lemma~12'', ``(38)'' or ``Table~1'' are
to its numbering. Our own equations and tables are numbered within
sections, as in (2.1) or Table~1.1, so a plain number always refers to
\OAI. Logarithms $\log$ are natural; $\lg$ is as in \OAI.

\begin{theorem}\label{thm:main}
Theorem~1 of \OAI{} holds with $\kappa=2^{-74}$. That is, there is one
deterministic Turing machine with a fixed finite alphabet and a fixed
finite number of one-dimensional tapes that computes the exact product
of two $n$-bit integers for every $n\ge1$ in worst-case time
$O\bigl(n(\lg n)^{1-2^{-74}}\bigr)$. Consequently Corollary~2 (exact
division and integer square root) and the transposition bound of
Section~11 of \OAI{} also hold with $\kappa=2^{-74}$.
\end{theorem}

The machine is the one constructed in \OAI, with five differences:
\begin{itemize}
\item the schedule of the axis movements in Lemma~20 and Section~9.3
(Section~\ref{sec:axes});
\item the size of the bit network (Section~\ref{sec:network});
\item the bound on the guard width in Section~5.4
(Section~\ref{sec:guard});
\item the Gaussian width $\alpha$ of Section~9.2.1, with the explicit
cutoff $b\ge2^{24}$ of Sections~9.2.1 and~9.7.1 raised
(Section~\ref{sec:alpha});
\item the fixed constants, including $\beta$, which Section~5.3 fixes
to $1/2$ (Section~\ref{sec:constants}).
\end{itemize}
Every other statement and proof of \OAI{} is used as it stands. Those
statements depend on the constants only through the inequalities
collected in Section~\ref{sec:constants}. We do not re-verify the parts
of \OAI{} that we do not change.

\paragraph{Where $2^{-182}$ comes from.}
Section~9.8 of \OAI{} divides the running time by $Tp=\Theta(n)$ and
bounds the quotient by seven powers $p^{1-g_i}$; any $\kappa<\min_i g_i$
is then admissible. With the constants (38), the smallest margin is
\[
 g_2=\epsilon c(1-\tau)=2^{-75}\cdot2^{-56}\cdot2^{-50}=2^{-181},
\]
a product of three small numbers.
\begin{itemize}
\item $1-\tau$ is the power saving of the chunk interchange
(Lemma~12). Section~3.6 of \OAI{} takes $1-\tau=2^{-50}$. The $h=100$
network actually gives $1.19\cdot2^{-40}$, and the network with $h=46$
gives $1.24\cdot2^{-36}$ (Section~\ref{sec:network}).
\item $c$ is the exponent of the chunk length $K=\lfloor d^c\rfloor$.
Proposition~14 requires $\tau(1+c/\beta)<\lambda<1$, hence
$c<\beta(1-\tau)/\tau$. With $\beta=1/2$ this forces $c$ below about
$2^{-51}$, and \OAI{} takes $c=2^{-56}$.
\item $\epsilon$ is the exponent of the number of axes
$d=\lfloor b^\epsilon\rfloor$. The margin $g_4=1-\tau-\epsilon(2-\tau)$
of the row ``CRT and axis layouts'' of Table~1 of \OAI{} must be positive, so
$\epsilon<(1-\tau)/(2-\tau)<2^{-50}$. This constraint comes only from the
$O(d^2)$ adjacent axis moves in Lemma~20 and Section~9.3.
\end{itemize}
Table~\ref{tab:ladder} records the effect of the changes, applied
cumulatively.

\begin{table}[ht]
\centering
\small
\renewcommand{\arraystretch}{1.3}
\begin{tabular}{@{}>{\raggedright\arraybackslash}p{.39\linewidth}lllll@{}}
\toprule
Change & $1-\tau$ & $c$ & $\epsilon$ & $\min_i g_i$ & $\kappa$\\
\midrule
\OAI, constants (38) & $2^{-50}$ & $2^{-56}$ & $2^{-75}$ & $2^{-181}$ & $2^{-182}$\\
$O(d)$ axis interchanges (Sec.~\ref{sec:axes}); new $\beta,c,\epsilon$
 & $2^{-50}$ & $\tfrac{31}{32}2^{-50}$ & $\tfrac1{15}$ & $2^{-103.95}$ & $2^{-104}$\\
admissible exponent of the $h=100$ bit network (Rem.~\ref{rem:networks}) & $2^{-40}$ & $\tfrac{31}{32}2^{-40}$ & $\tfrac1{15}$ & $2^{-83.95}$ & $2^{-84}$\\
bit network with $h=46$ (Sec.~\ref{sec:network})
 & $2^{-36}$ & $\tfrac{31}{32}2^{-36}$ & $\tfrac1{15}$ & $2^{-75.95}$ & $2^{-76}$\\
guard width $C_1=5$, smaller $\alpha$, $\delta=\tfrac1{128}$, and $1-\tau=\tfrac65 2^{-36}$ (Secs.~\ref{sec:network}--\ref{sec:alpha})
 & $\tfrac65 2^{-36}$ & $\tfrac76 2^{-36}$ & $\tfrac3{16}$ & $2^{-73.93}$ & $2^{-74}$\\
\bottomrule
\end{tabular}
\caption{Cumulative effect of the changes. Rows 2--4 use
$C_1=14$ (which \OAI{} notes is sufficient; it takes $C_1=20$), the Gaussian width
$\alpha$ of \OAI{} and $\delta=1/16$. All rows after the first use
$\beta=1-2^{-10}$, $\sigma=1-2^{-50}$ and $\lambda'=1-2^{-61}$. Row~2 uses
$\lambda=1-2^{-56}$ and the others $\lambda=1-2^{-51}$. In every row
the minimum margin is $g_2$, and $\kappa$ is the largest power of two
below it.}
\label{tab:ladder}
\end{table}

Section~\ref{sec:limits} shows that the coupling between $c$ and $\tau$
makes $\kappa\approx\epsilon(1-\tau)^2$ unavoidable in this
architecture. Section~\ref{sec:verification} describes the
verification.

\section{Moving axes with \texorpdfstring{$O(d)$}{O(d)} interchanges}
\label{sec:axes}

\begin{lemma}[Interchanging two axis fields]\label{lem:direct-swap}
Consider a bit array of volume $V$ whose address consists of fields
$F_1,\dots,F_k$, in this physical order, followed by a record suffix.
Let $i<j$ and suppose $F_i,F_j$ have widths in $\{\ell-1,\ell\}$. The
address permutation that exchanges the contents of $F_i$ and $F_j$, so
that the old $F_j$ occupies position $i$ and the old $F_i$ position $j$,
and fixes all other fields, can be performed in $O(V(1+\ell^\tau))$
steps on a fixed number of tapes. The bound includes descriptor work.
\end{lemma}

\begin{proof}
Collapse the fields before $F_i$ into a prefix $[P]$, the fields
strictly between $F_i$ and $F_j$ into a gap $[G]$, and the fields after
$F_j$ together with the record suffix into $[B]$. Lemma~12 of \OAI{}
accepts arbitrary positive $P,G,B$. If the two widths are equal, it
performs the interchange in $O(V\ell^\tau)$ steps. If they differ by
one bit, apply the remark following Lemma~12: write the longer field as
a leading bit followed by a field of the shorter width. One equal-width
interchange (Lemma~12) and one single-bit move (Lemma~3, $O(V)$) give
the required permutation, and the gap stays between the two full
fields. The collapsed lengths have $O(\log V)$ bits, and computing them
takes $(\log 2V)^{O(1)}=O(V)$ steps.
\end{proof}

\begin{lemma}[Replacement for the movement bound in Lemma~20]
\label{lem:tensor-moves}
Under the hypotheses of the second paragraph of Lemma~20 of \OAI{}
($t_i\in\{r/2,r\}$, the full padded box of $T$ entries, $O(p)$ bits per
entry), exposing and restoring all the lines of the linewise
implementations of $\mathcal A$ and $\mathcal B_0$ costs
\[
 O\bigl(Tp\,d\,(1+\ell^\tau)\bigr).
\]
\end{lemma}

\begin{proof}
In the proof of Lemma~20, replace the schedule of adjacent interchanges
by the following. To expose axis $i$, interchange its field with the
innermost axis field by Lemma~\ref{lem:direct-swap}, unless axis $i$ is
already innermost. Then process its lines exactly as in \OAI{} and
interchange the two fields back.

The descriptor records which coordinate each physical field holds.
Moving fields therefore changes their positions but not the validity
conditions, as already noted in the proof of Lemma~20. The line maps
write $t_i$ values per line, so the padded box keeps its shape and the
field widths never change. After axis~$i$ is processed the physical
order is the original one. Hence later steps, which expect axis-major
order with the length-$r$ axis last, are unaffected.

Each tensor map performs at most $2d$ interchanges, each costing
$O(Tp(1+\ell^\tau))$, and there is a fixed number of tensor maps. The
interchanges are exact permutations of the padded box, so they change
neither norms nor errors.
\end{proof}

\begin{lemma}[Replacement in Section~9.3]\label{lem:crt-moves}
The coefficient-address map $\Phi$ of Section~9.3 of \OAI, its inverse
and their padding operations cost $O(Tp\,d\,(1+\ell^\tau))$.
\end{lemma}

\begin{proof}
The only change concerns the step that reorders the complete padded
axes to $(a_1,\ldots,a_d)$. The nested insertion scan produces the
physical order $(a_d,\ldots,a_1)$. Reversing it takes the $\lfloor
d/2\rfloor$ interchanges of the fields in positions $k$ and $d+1-k$,
$1\le k\le\lfloor d/2\rfloor$, each by Lemma~\ref{lem:direct-swap}
(widths $\ell$ or $\ell-1$).

All of these interchanges happen before any triangular update. So when
target $i$ is updated, in the order $i=d,\dots,1$, its lower-index
controls precede it, as Lemma~4 requires. The inverse first performs the
inverse triangular updates ($i=1,\dots,d$) and then the same $\lfloor
d/2\rfloor$ interchanges. The $d$ controlled rotations still cost
$O(dV)$, and their offset preparation is $o(Tp)$, as in \OAI.
\end{proof}

Consequently the row ``CRT and axis layouts'' of Table~1 of \OAI{} becomes
$d(1+\ell^\tau)$. Since $d=\Theta(p^\epsilon)$ and
$\ell=\Theta(p^{1-\epsilon})$, its power of $p$ is
$\epsilon+\tau(1-\epsilon)$ and its margin is
\begin{equation}\label{eq:g4}
 g_4=(1-\tau)(1-\epsilon),
\end{equation}
which is positive for every $0<\epsilon<1$. The condition
$\epsilon(2-\tau)<1-\tau$ of Section~9.1 disappears. The binding upper
bounds on $\epsilon$ come from the guard width ($\epsilon C_1<1$) and
the Gaussian row (Sections~\ref{sec:guard} and~\ref{sec:alpha}).

\section{A smaller address-interchange network}\label{sec:network}

The bit network of Section~3.1 of \OAI{} is defined for every ground-set
size~$h$. \OAI{} uses $h=100$ for both networks, but notes in
Section~5.3 that the constants of the complex network ``need not equal
the constants used to prove the chunk-swap lemma''. We keep the complex
(phase) network at $h=100$, with $\sigma=1-2^{-50}$ as in~(10), and use
the bit network with $h=46$.

For general $h$ the bit construction has
\[
 v=\tbinom h3,\quad N=v^3,\quad m=h^3,\quad I=3v^2,\quad
 \zb=3\tbinom{h-3}{2},\quad \Wb=2N+I(v\zb+h),\quad \Lb=Ih^2,
\]
and $\sbb=\Wb m-N+2\Lb$ as in Proposition~7. Nothing in Sections~3.1--3.4
of \OAI{} depends on $h=100$ except the printed numbers. The form
$x^{\mathsf T}(I-J/9)y$ is nondegenerate whenever $h\ne9$ (its
eigenvalues are $1$ and $1-h/9$). Every label in table~(6) and in the
proof of Lemma~6 is built from lines spanned by vectors of nonzero norm
(the triple indicators satisfy $\langle t_T,t_T\rangle=2$, and
neighboring indicators are orthogonal), using orthogonal complements,
tensor products and orthogonal sums. So all labels are nondegenerate,
and Lemma~6 and Proposition~7 hold with these counts. The binary
orthonormal-basis assertion of Lemma~6 concerns only the complex
network. The rank saving is $\Wb m-\sbb=N-2\Lb=v^2(v-6h^2)$, which is
positive if and only if $h\ge39$.

For $h=46$:
\[
\begin{aligned}
 &v=15180,\qquad N=3\,497\,963\,832\,000,\qquad m=97\,336,\qquad
  I=691\,297\,200,\qquad \zb=2709,\\
 &\Wb=28\,434\,979\,789\,999\,200,\qquad
  \Lb=1\,462\,784\,875\,200=\tfrac{23}{55}N,\\
 &\sbb=2\,767\,747\,192\,266\,968\,049\,600,\qquad
  \Wb m=2\,767\,747\,192\,839\,362\,131\,200,\\
 &\etab:=\frac{\Wb m-\sbb}{\Wb m}=\frac{9}{43\,518\,487\,588}
  =2.0681\ldots\cdot10^{-10}.
\end{aligned}
\]

\begin{lemma}\label{lem:tau}
For the bit network with $h=46$ and $\tau=1-\tfrac65 2^{-36}$, and
\emph{a fortiori} for $\tau=1-2^{-36}$, we have $\sbb/\Wb<m^\tau$. So
Proposition~11 and Lemma~12 hold with this $\tau$.
\end{lemma}

\begin{proof}
Put $T=1-\tau$. Since $m<2^{17}$ and $\log2<0.6932$, we have $\log
m<11.7844$, and therefore
\[
 T\log m<\tfrac65\cdot11.7844\cdot2^{-36}<2.058\cdot10^{-10}<\etab.
\]
Hence $m^\tau=m\exp(-T\log m)>m(1-T\log m)>m(1-\etab)=\sbb/\Wb$.
\end{proof}

\begin{remark}\label{rem:networks}
We record three facts about the two networks.
\begin{enumerate}
\item The exact exponent is
$\log_m(\sbb/\Wb)=1-T^*(46)$ with $T^*(46)=1.23732\ldots\cdot2^{-36}$.
Over all ground-set sizes, $T^*(h):=-\log(1-\etab(h))/\log h^3$ is
largest at $h=46$; see Table~\ref{tab:h} in Appendix~\ref{app:h}.
\item Even the $h=100$ bit network of \OAI{} has
$T^*(100)=1.194\ldots\cdot2^{-40}$, so $\tau=1-2^{-40}$ is admissible
there. Section~3.6 of \OAI{} uses only $1-2^{-50}$.
\item Lemma~10, Proposition~11 and Lemma~12 are generic in
$(m,W,s)$. Only the finite rational table of Lemma~10 and its prime
$q$ must be recomputed for the matrices of the new network. In
Section~5 of \OAI, $\tau$ enters only through the cost $O(V(L^\tau+1))$ of
Lemma~13, that is, through the conditions $\max\{\tau,\sigma\}<\lambda$
and $\tau(1+c/\beta)<\lambda$ of Proposition~14.
\end{enumerate}
\end{remark}

\section{The guard width}\label{sec:guard}

Section~5.4 of \OAI{} bounds the arithmetic depth of one simultaneous
layer by the recursion
\begin{equation}\label{eq:depth-rec}
 A(e)\le s\,A(e/m)+E\quad\text{(internal node)},\qquad
 A(e)\le 8e\quad\text{(leaf, $e<d^\beta$)},
\end{equation}
where $m$, $s=\scc$ and $W=\Wc$ are the constants of the complex network
($h=100$) and $E=64(W+m+1)^3$. It then uses $sx+E\le(s+E)x=Bx$, which
replaces the per-level factor $s\approx m^{4.55}$ by $B\approx m^{10.94}$
and leads to the exponent $\nu+3=14$ with $\nu=11$ (\OAI{} takes
$C_1=20$). Unrolling the recursion additively loses much less.

\begin{lemma}\label{lem:unroll}
Let $M\ge1$ and $C\ge1$ be real numbers, let $E\ge0$, and let
$2\le\mu\le M^C$.
Suppose $A_0,A_1,\ldots$ satisfy $A_j\le 8M^j$ at leaf levels (with $0$
a leaf level) and $A_j\le\mu A_{j-1}+E$ at all other levels. Then
$A_j+E\le(8+E)M^{Cj}$ for all $j$.
\end{lemma}

\begin{proof}
Induction on $j$. At a leaf level, $A_j+E\le 8M^j+E\le(8+E)M^{Cj}$. At
any other level, $A_j+E\le\mu A_{j-1}+2E\le\mu(A_{j-1}+E)\le
M^C(8+E)M^{C(j-1)}$, using $\mu\ge2$.
\end{proof}

\begin{corollary}[Guard exponent $C_1=5$]\label{cor:C1}
Section~5.4 and Proposition~14 of \OAI{} hold with $C_1=5$, that is,
with $\Delta=\lceil C_0d^5\rceil$ and the same $C_0=128mB^2$. The only
resulting condition is $\epsilon C_1<1$.
\end{corollary}

\begin{proof}
For the $h=100$ complex network,
$2\le\scc=1\,873\,807\,244\,636\,671\,267\,308\,000\,000<10^{30}=m^5$.
By Lemma~\ref{lem:unroll}, applied to \eqref{eq:depth-rec} with
$A_j=A(m^j)$, a piece of $m^j$ chunks has depth at most
$(8+E)m^{5j}$. Write $D-q_0=\sum_ja_jm^j$ as in \OAI. Since $\sum
x_i^5\le(\sum x_i)^5$, the pieces together contribute at most
$(8+E)(D-q_0)^5\le(8+E)d^5$. At most $d$ individually processed kernels
add $8d$, and the outer diagonals and the scalar $b_D$ add $t+9\le
d+9$. The depth of a complete layer is therefore at most
$(26+E)d^5\le C_0d^5$, because $C_0=128mB^2\ge26+E$. The remainder of
Section~5.4 uses $C_1$ only through $\Delta\le p$, that is,
$\epsilon C_1<1$.
\end{proof}

\begin{remark}[Guard exponent $1+2^{-10}$]\label{rem:tv}
The recursion~\eqref{eq:depth-rec} concatenates all $s$ child calls of
a network invocation, but a dependency chain meets far fewer.

A chain meets only the child calls on the edges of the wire it
currently occupies, and it changes wires only at gates, where all
incident wires carry the same label. So the number of child calls a
chain meets at one level is the total variation of the label dimension
along a path in the label graph given by table~(6) and Lemma~6.

Let $k(x)$ count the decreasing center edges ($3\to4$, one per stage)
that can precede the vertex $x$, and put
$\Phi(x)=\dim x+2h\,k(x)$. Along every edge, $\Phi$ increases by at
least the absolute change in dimension, and $0\le\Phi\le m+6h$. Hence every chain
meets at most $m+6h$ child calls per level, and this bound is attained.

Lemma~\ref{lem:unroll} with $\mu=m+6h\le m^{1+2^{-10}}$ then gives the
guard exponent $1+2^{-10}$, and with $\mu\le m^2$ it gives $C_1=2$. We
do not need this sharper bound.
\end{remark}

\section{A smaller Gaussian width}\label{sec:alpha}

Lemma~19 of \OAI{}, which adapts the Gaussian resampling of
Harvey and van der Hoeven~\cite[Theorem~4.2, Proposition~4.7 and
Lemmas~4.8--4.12]{HvdH},
requires $2\le\alpha<\sqrt p$ and $\theta>p/\alpha^4$,
and in Section~9.2.2 every $\theta_i=t_i/s_i-1$ exceeds
$\eta=1/(4d)$. Hence
\begin{equation}\label{eq:alpha}
 \alpha:=\bigl\lceil(24db)^{1/4}\bigr\rceil
\end{equation}
suffices, since $\alpha^4\theta_i>24db/(4d)=6b=p$; and $2\le\alpha<\sqrt
p$ for all large $n$. \OAI{} takes $\alpha=\lceil(12d^2b)^{1/4}\rceil$,
which carries an extra factor $d/2$ inside the fourth root. With
\eqref{eq:alpha}:
\begin{itemize}
\item $\alpha=\Theta(p^{(1+\epsilon)/4})$, so the row ``Gaussian line
maps'' $dp^{1/2+\delta}\alpha$ of Table~1 of \OAI{} has power
$3/4+\delta+5\epsilon/4$ and margin
$g_5=1/4-\delta-5\epsilon/4$;
\item $\gamma=2d\alpha^2\le8\sqrt{24}\,d^{3/2}b^{1/2}<40\,b^{1/2+3\epsilon/2}$.
This is $o(b)$ for $\epsilon<1/3$, so (44) holds for all large $n$. The
explicit cutoff $b\ge2^{24}$ in Sections~9.2.1 and~9.7.1 must be raised;
for $\epsilon=3/16$, $b\ge2^{34}$ gives $\gamma\le b/4$;
\item the Neumann count $k\le\alpha^2+1<p+1$ and the cost bound
$O(tp^{3/2+\delta}\alpha)$ of Lemma~19 are unchanged, since their
proofs use only $\alpha<\sqrt p$ and $\alpha^4\theta>p$.
\end{itemize}

\section{The new constants}\label{sec:constants}

Replace (38) of \OAI{} by
\begin{equation}\label{eq:new-constants}
\begin{gathered}
 \tau=1-\tfrac65 2^{-36}\ \ (\text{bit network, }h=46),\qquad
 \sigma=1-2^{-50}\ \ (\text{complex network, }h=100),\\
 \beta=1-2^{-10},\qquad c=\tfrac76 2^{-36},\qquad
 \lambda=1-2^{-51},\qquad \lambda'=1-2^{-61},\\
 C_1=5,\qquad \epsilon=\tfrac3{16},\qquad \delta=\tfrac1{128},\qquad
 \kappa=2^{-74}.
\end{gathered}
\end{equation}
Section~5.3 and the statement of Proposition~14 of \OAI{} fix
$\tau=\sigma=1-2^{-50}$, $\beta=1/2$ and $C_1=20$. The proofs are
parametric in these constants (for $C_1$, use
Corollary~\ref{cor:C1}), and we read Proposition~14 with the values
\eqref{eq:new-constants}. Its conditions hold:
\[
 \max\{\tau,\sigma\}=1-2^{-50}<\lambda<1,\qquad
 \tau(1+c/\beta)<1-2^{-41}<\lambda,
\]
\[
 \sigma+\beta(1-\sigma)=1-2^{-60}<\lambda'<1,\qquad
 \lambda<\lambda',\qquad
 \epsilon C_1=\tfrac{15}{16}<1.
\]
The remaining conditions of Section~9 also hold:
\begin{itemize}
\item $0<\delta<1/8$ (Lemma~19);
\item $\epsilon<1/3$ ($\gamma=o(b)$, Section~\ref{sec:alpha});
\item $\epsilon<1/2$ (prime search: $\log(r/2)/d=\Theta(p^{1-2\epsilon})\to\infty$);
\item $\epsilon c>0$ ($K/\log p\to\infty$);
\item $K=\lfloor d^c\rfloor\le\ell-1$ for large $n$.
\end{itemize}

\begin{table}[ht]
\centering
\small
\begin{tabular}{@{}>{\raggedright\arraybackslash}p{.42\linewidth}ll@{}}
\toprule
Operation & Normalized cost & Power of $p$\\
\midrule
Prefix-slot moves and individual butterfly rounds & $dK$ & $\epsilon(1+c)$\\
Chunk exchanges for transform layout & $pK^{\tau-1}$ & $1-\epsilon c(1-\tau)$\\
Simultaneous butterfly rounds & $\ell d^{\lambda'}$ & $1-\epsilon(1-\lambda')$\\
CRT and axis layouts (Sec.~\ref{sec:axes}) & $d(1+\ell^\tau)$ & $1-(1-\tau)(1-\epsilon)$\\
Gaussian line maps (Sec.~\ref{sec:alpha}) & $dp^{1/2+\delta}\alpha$ & $3/4+\delta+5\epsilon/4$\\
Chirps, twists and scalar products & $dp^\delta$ & $\epsilon+\delta$\\
Packed polynomial products & $\log(rp)$ & $1-\epsilon$\\
Final scaling and rounding & $p^\delta$ & absorbed above\\
Input, padding, carries and output & $1$ & absorbed above\\
\bottomrule
\end{tabular}
\caption{Replacement for Table~1 of \OAI: costs divided by $Tp$.}
\label{tab:costs}
\end{table}

The margins of the seven dominant rows of Table~\ref{tab:costs} are
\begin{equation}\label{eq:margins}
\begin{aligned}
 g_1&=1-\epsilon(1+c)=\tfrac{13}{16}-\tfrac7{32}2^{-36}, &
 g_2&=\epsilon c(1-\tau)=\tfrac{21}{80}\,2^{-72}=1.05\cdot2^{-74},\\
 g_3&=\epsilon(1-\lambda')=\tfrac3{16}2^{-61}, &
 g_4&=(1-\tau)(1-\epsilon)=\tfrac{39}{40}2^{-36},\\
 g_5&=\tfrac14-\delta-\tfrac54\epsilon=\tfrac1{128}, &
 g_6&=1-\delta-\epsilon=\tfrac{103}{128},\qquad g_7=\epsilon=\tfrac3{16}.
\end{aligned}
\end{equation}
Thus $\min_ig_i=g_2>2^{-74}=\kappa$.

\begin{proof}[Proof of Theorem~\ref{thm:main}]
Follow Section~9.8 of \OAI{} with Table~\ref{tab:costs} in place of
Table~1 of \OAI. Each row is $O(p^{1-g_i})$: the factor $\log(2d)$ from the
base-$m$ groups is absorbed in $\lambda'$, and $\log p$ in $p^\delta$.
Since $\kappa<\min_ig_i$, any fixed power of $\log p$ is at most
$p^{\min_ig_i-\kappa}$ for large $p$. The factor two in ``$\min_i
g_i=2\kappa$'' of \OAI{} is not needed, although here it would cost only
one bit.

Hence the large-input cost is
$O(Tp\cdot p^{1-\kappa})=O(n(\lg n)^{1-\kappa})$. The finitely many
small lengths are handled by schoolbook multiplication, as in \OAI.
Corollary~2 and Section~11 use only the multiplication bound and
$0<\kappa<1$.
\end{proof}

\begin{remark}[$\kappa=2^{-76}$ from Sections~\ref{sec:axes}
and~\ref{sec:network} alone]\label{rem:setA}
Use the guard exponent $C_1=14$, which \OAI{} notes is sufficient (it
takes $C_1=20$); keep its Gaussian width $\alpha$ and $\delta=1/16$. Take
$\tau=1-2^{-36}$ ($h=46$), $c=\tfrac{31}{32}2^{-36}$ and $\epsilon=1/15$,
with $\sigma,\beta,\lambda,\lambda'$ as in \eqref{eq:new-constants}. All
conditions hold, including $\epsilon<1/12$ as used in~(40). The margins
have minimum $g_2=\tfrac{31}{480}2^{-72}>2^{-76}$, with
$g_5=1/4-\delta-3\epsilon/2=7/80$, so $\kappa=2^{-76}$.
\end{remark}

\section{Limits of parameter tuning}\label{sec:limits}

\begin{proposition}\label{prop:barrier}
Keep the algorithm of \OAI{} with the changes above and its cost
accounting (Table~\ref{tab:costs}). Then every admissible choice of
constants satisfies $\kappa<\epsilon(1-\tau)^2/\tau$ and
$\epsilon<1/5$. If the interchange network is the bit network of
Section~3.1 of \OAI{} for some ground-set size~$h$, then also
$1-\tau<1.2374\cdot2^{-36}$, and consequently $\kappa<2^{-73.70}$.
\end{proposition}

\begin{proof}
The condition $\tau(1+c/\beta)<\lambda<1$ of Proposition~14 gives
$c<\beta(1/\tau-1)<(1-\tau)/\tau$, where $\beta<1$ because
$\sigma+\beta(1-\sigma)<\lambda'<1$; so $\kappa<g_2=\epsilon
c(1-\tau)<\epsilon(1-\tau)^2/\tau$. The condition $g_5>0$ gives
$\epsilon<1/5-4\delta/5<1/5$; with the Gaussian width of \OAI{} it gives
even $\epsilon<1/6$.

Proposition~11 needs $s/W<m^\tau$, that is, $1-\tau<T^*(h)$. For $h<39$
there is no rank saving. For $39\le h\le2000$ a numerical computation
(Appendix~\ref{app:script}; selected values in Table~\ref{tab:h}) gives $T^*(h)\le T^*(46)<1.2374\cdot2^{-36}$. For
$h>2000$,
\[
 \etab\le\frac{N}{\Wb m}\le\frac{N}{Iv\zb m}=\frac1{3\zb m},
\]
and therefore $T^*(h)<2\etab/\log m<10^{-18}$. Finally
$\tfrac15(1.2374\cdot2^{-36})^2/(1-1.2374\cdot2^{-36})<2^{-73.70}$.
\end{proof}

The obstruction is structural. The chunk length $K=\lfloor d^c\rfloor$
serves two purposes. In the transform layout of Lemma~16 there are
$O(p/K)$ chunk interchanges of width $K$, with normalized cost
$pK^{\tau-1}$, and this term wants $K$ large. In each simultaneous
layer, the packed selected-bit additions of Lemma~13 interchange slots
of width $(e/m)K$, at cost $O((eK)^\tau)$ per node, and this forces
$c<(1-\tau)/\tau$. Going substantially below $2^{-74}$ therefore needs a
new idea. Two possibilities are an address-interchange network with a
larger rank saving (the saving enters squared), or a layout whose chunk
length is not tied to the swap exponent, for instance a way to realize
the layout permutation of Lemma~16 at cost below $pK^{\tau-1}$.

\section{Verification}\label{sec:verification}

\paragraph{Exact arithmetic.} The script in Appendix~\ref{app:script}
(Python) checks every condition with integers and fractions; floating
point is used only for displayed logarithms and for the scan over $h$.
It does the following:
\begin{itemize}
\item recomputes the printed constants of \OAI{} for $h=100$;
\item computes the $h=46$ counts and $\etab$;
\item checks the exponent certificates of Lemma~\ref{lem:tau} and of
Remark~\ref{rem:networks}, and the optimality of $h=46$ for $h\le2000$;
\item checks every condition and margin for the five rows of
Table~\ref{tab:ladder}, including $\epsilon<1/12$ for the four rows that
keep the Gaussian width of \OAI;
\item checks $8\sqrt{24}<40$ and that $40\,b^{25/32}\le b/4$ for
$b\ge2^{34}$ (Section~\ref{sec:alpha});
\item evaluates the barrier of Proposition~\ref{prop:barrier}.
\end{itemize}

\paragraph{Lean.} The companion file \lean{IntMulKappa.lean} checks the
same arithmetic independently in Lean~4~\cite{Lean}. It uses core Lean only (no
Mathlib), version~4.34.1, the toolchain pinned by the
\texttt{openai/math} repository. It contains no \lean{sorry} and no
\lean{native\_decide}, and \lean{\#print axioms} reports at most
\lean{propext}, \lean{Classical.choice} and \lean{Quot.sound}. It
compiles in a few seconds with \lean{lean IntMulKappa.lean}. Its
main statements are listed in Table~\ref{tab:lean}.

\begin{table}[ht]
\centering
\small
\begin{tabular}{@{}>{\raggedright\arraybackslash}p{.27\linewidth}>{\raggedright\arraybackslash}p{.66\linewidth}@{}}
\toprule
Lean theorem & Content\\
\midrule
\lean{printed\_values\_h100}, \lean{counts\_h46}, \lean{relations\_h46},
\lean{eta\_b\_h46} & the constants of Section~\ref{sec:network} and of
\OAI{} (3)--(4), (9), Proposition~8\\
\lean{rate\_cert} & if $Q,Y>0$ and $Q^qM^p<y_n^q$, where
$y_0=\lfloor QX/Y\rfloor$ and $y_{i+1}=\lfloor y_i^2/Q\rfloor$, then
$Y^{q2^n}M^p<X^{q2^n}$\\
\lean{tau\_B}, \lean{tau\_A} & $\sbb^{K}m^{6}<(\Wb m)^{K}$ with
$K=5\cdot2^{36}$, and $\sbb^{2^{36}}m<(\Wb m)^{2^{36}}$ ($h=46$); these
are equivalent to $\sbb/\Wb<m^{1-\frac65 2^{-36}}$ and
$\sbb/\Wb<m^{1-2^{-36}}$\\
\lean{tau\_100}, \lean{sigma\_paper} & the same for the $h=100$ bit
network with $1-2^{-40}$, and for the $h=100$ complex network with
$\sigma=1-2^{-50}$\\
\lean{depth\_unroll}, \lean{sum\_pow\_le} & Lemma~\ref{lem:unroll} for integer $C$,
and $\sum x_i^C\le(\sum x_i)^C$\\
\lean{guard\_instances\_h100} & $2\le\scc\le m^5$ and $m+6h\le m^2$ for $h=100$\\
\lean{chain\_tv\_bound} ($h=100$), \lean{chain\_tv\_bound\_25},
\lean{chain\_tv\_bound\_46}, \lean{extremal\_tv} & in the label graph of
Remark~\ref{rem:tv}, every path has total variation at most $h^3+6h$,
and one path attains it\\
\lean{SetB.certified}, \lean{SetA.certified}, \lean{PaperSet.certified}
& all conditions and margins of \eqref{eq:new-constants}, of
Remark~\ref{rem:setA}, and of (38) of \OAI{} ($\min g=g_2=2^{-181}$)\\
\bottomrule
\end{tabular}
\caption{Main statements of \lean{IntMulKappa.lean}.}
\label{tab:lean}
\end{table}

The exponents $q2^n$ in \lean{tau\_B} and the other certificates are far
too large to evaluate. These statements are proved by applying
\lean{rate\_cert}, and the kernel only checks the small final
inequality, which needs 36 to 50 truncated squarings at scale
$Q=2^{160}$.

The Lean file certifies the exact arithmetic of the improvement: the
network counts, the exponent certificates, and all conditions and
margins. The scan over $h$ and the barrier constant of
Proposition~\ref{prop:barrier} are numerical (Appendix~\ref{app:script}).
The file does not formalize the machine model, the algorithm of \OAI, or
Lemmas~\ref{lem:direct-swap}--\ref{lem:crt-moves}, and the
correspondence between dependency chains and paths in the label graph is
part of the informal argument.

At the time of writing, the \texttt{openai/math} repository
(commit \texttt{adc7f12}) contains no formalization of \OAI{}
(family~109). Its formalized Fourier result,
\lean{OAI.ExactFourier.main\_theorem} (family~130), is a different
statement in a different model: a subsequential gate-count bound for
exact complex Fourier circuits (for every $c>0$ there are arbitrarily
large $n$ with a DFT circuit of fewer than $cn\log_2n$ gates).

\appendix

\section{Exponents of the bit networks}\label{app:h}

\begin{table}[H]
\centering
\small
\begin{tabular}{@{}rll@{\qquad}rll@{}}
\toprule
$h$ & $\etab(h)$ & $T^*(h)/2^{-36}$ & $h$ & $\etab(h)$ & $T^*(h)/2^{-36}$\\
\midrule
39 & $4.228\cdot10^{-12}$ & 0.0264 & 47 & $2.065\cdot10^{-10}$ & 1.2288\\
40 & $7.385\cdot10^{-11}$ & 0.4586 & 48 & $2.037\cdot10^{-10}$ & 1.2052\\
41 & $1.234\cdot10^{-10}$ & 0.7614 & 49 & $1.990\cdot10^{-10}$ & 1.1710\\
42 & $1.579\cdot10^{-10}$ & 0.9678 & 50 & $1.929\cdot10^{-10}$ & 1.1298\\
43 & $1.810\cdot10^{-10}$ & 1.1023 & 51 & $1.861\cdot10^{-10}$ & 1.0840\\
44 & $1.955\cdot10^{-10}$ & 1.1833 & 52 & $1.786\cdot10^{-10}$ & 1.0356\\
45 & $2.035\cdot10^{-10}$ & 1.2248 & 100 & $1.501\cdot10^{-11}$ & 0.0747\\
46 & $2.068\cdot10^{-10}$ & \textbf{1.2373} & & &\\
\bottomrule
\end{tabular}
\caption{Rank-saving ratio $\etab(h)=(N-2\Lb)/(\Wb m)$ and exact swap
exponent $1-T^*(h)$, $T^*(h)=-\log(1-\etab)/\log h^3$, of the bit
network of \OAI, Section~3.1, near the optimum. For $h\le38$ there is no
saving. The value at $h=100$ is $T^*(100)=1.194\cdot2^{-40}$.}
\label{tab:h}
\end{table}

\section{Verification script}\label{app:script}

\begin{lstlisting}[language=Python]
"""Verification of the constants in the note "Integer multiplication below n log n:
the exponent kappa = 2^-74".  All conditions and margins use exact integer/Fraction
arithmetic; floating-point logarithms are used only for display and for the scan of
T*(h) over h.  Run: python3 verify_kappa.py"""
from fractions import Fraction as F
from math import comb, log, log2

def bit_net(h):          # Section 3.1 of [OAI], bit construction, ground set [h]
    v = comb(h, 3); N = v**3; m = h**3; I = 3*v*v
    W = 2*N + I*(v*3*comb(h-3, 2) + h); L = I*h*h
    return dict(v=v, N=N, m=m, W=W, L=L, s=W*m - N + 2*L)
def cx_net(h):           # Section 3.1 of [OAI], complex construction
    v = comb(h, 3); N = v**3; m = h**3; I = 3*v*v
    W = 2*N + I*(v*(comb(h-3, 3) + 3*(h-3)) + h + 1); L = I*(h+1)*h
    return dict(v=v, N=N, m=m, W=W, L=L, s=W*m - 2*N + 2*L)
def eta(d): return F(d['W']*d['m'] - d['s'], d['W']*d['m'])
def Tstar(d):            # 1 - log_m(s/W) = -ln(1-eta)/ln m  (floating point, for display)
    e = eta(d); return -log(1 - e.numerator/e.denominator)/log(d['m']) if e > 0 else float('nan')

# --- the paper's printed values (h = 100) -------------------------------------------------
b100, c100 = bit_net(100), cx_net(100)
assert (b100['W'], b100['s']) == (177176569091445000000, 177176569088785861287000000)
assert (c100['W'], c100['s']) == (1873807244643542670000, 1873807244636671267308000000)
assert eta(b100) == F(339, 22587335000000) and eta(c100) == F(73, 19906842167500)

# --- the h = 46 bit network ---------------------------------------------------------------
b = bit_net(46)
assert (b['v'], b['N'], b['m']) == (15180, 3497963832000, 97336)
assert (b['W'], b['L'], b['s']) == (28434979789999200, 1462784875200, 2767747192266968049600)
assert F(b['L'], b['N']) == F(23, 55) and eta(b) == F(9, 43518487588)
print("h=46: W_b m =", b['W']*b['m'], "  eta_b =", eta(b), " T* = %.5f * 2^-36" % (Tstar(b)*2**36))

# exponent certificates:  T ln m < eta  ==>  m^(1-T) > m(1 - T ln m) > m(1-eta) = s/W
ln2_up = F(6932, 10000); assert log(2) < 0.6932
def cert(T, d, ln_m_up): return T*ln_m_up < eta(d)
assert b['m'] < 2**17 and cert(F(6, 5)/2**36, b, 17*ln2_up)          # tau = 1-(6/5)2^-36
assert cert(F(1, 2**36), b, 17*ln2_up)                                # tau = 1-2^-36
assert b100['m'] < 2**20 and cert(F(1, 2**40), b100, 20*ln2_up)       # h=100: tau = 1-2^-40
assert cert(F(1, 2**50), c100, 20*ln2_up)                             # sigma = 1-2^-50 (paper)

# optimality of h = 46 among the bit networks (deficit N - 2L > 0 needs h >= 39)
assert all(bit_net(h)['N'] <= 2*bit_net(h)['L'] for h in range(7, 39))
best = max(range(39, 2001), key=lambda h: Tstar(bit_net(h)))
assert best == 46
print("T*(h) for h = 39..52, 100:", {h: round(Tstar(bit_net(h))*2**36, 4) for h in list(range(39, 53)) + [100]})

# guard width: s_c <= m^5 (C1 = 5) and the chain-variation route m + 6h <= m^(1+2^-10)
assert 2 <= c100['s'] <= c100['m']**5 and (1 + 600/10**6)**1024 <= 10**6

# --- parameter sets ------------------------------------------------------------------------
def check(name, T, S, beta, c, lam, lamp, C1, eps, delta, a5, kappa, axis_linear=True):
    tau, sig = 1 - T, 1 - S
    eps_max = F(1, 12) if a5 == F(3, 2) else F(1, 3)   # (40) of [OAI] / gamma = o(b) with the new alpha
    conds = [max(tau, sig) < lam < 1, tau*(1 + c/beta) < lam, max(lam, sig + beta*(1 - sig)) < lamp < 1,
             eps*C1 < 1, 0 < delta < F(1, 8), 0 < c, 0 < beta < 1, eps < eps_max]
    g = [1 - eps*(1 + c), eps*c*T, eps*(1 - lamp),
         T*(1 - eps) if axis_linear else 1 - tau - eps*(2 - tau),
         F(1, 4) - delta - a5*eps, 1 - delta - eps, eps]
    assert all(conds), name
    assert all(x > 0 for x in g) and kappa < min(g), name
    print(f"{name:7s} min g = 2^{log2(float(min(g))):.4f}  (g2 = {g[1]*2**72} * 2^-72),  kappa = 2^{log2(float(kappa)):.0f}")
    return min(g)
check("paper", F(1, 2**50), F(1, 2**50), F(1, 2), F(1, 2**56), 1 - F(1, 2**52), 1 - F(1, 2**54),
      20, F(1, 2**75), F(1, 16), F(3, 2), F(1, 2**182), axis_linear=False)
beta, lam, lamp, S = 1 - F(1, 2**10), 1 - F(1, 2**51), 1 - F(1, 2**61), F(1, 2**50)
check("step1", F(1, 2**50), S, beta, F(31, 32)/2**50, 1 - F(1, 2**56), lamp, 14, F(1, 15), F(1, 16), F(3, 2), F(1, 2**104))
check("step2", F(1, 2**40), S, beta, F(31, 32)/2**40, lam, lamp, 14, F(1, 15), F(1, 16), F(3, 2), F(1, 2**84))
check("Set A", F(1, 2**36), S, beta, F(31, 32)/2**36, lam, lamp, 14, F(1, 15), F(1, 16), F(3, 2), F(1, 2**76))
gB = check("Set B", F(6, 5)/2**36, S, beta, F(7, 6)/2**36, lam, lamp, 5, F(3, 16), F(1, 128), F(5, 4), F(1, 2**74))
assert gB == F(21, 80)/2**72
TB, cB = F(6, 5)/2**36, F(7, 6)/2**36
assert (1 - TB)*(1 + cB/beta) < 1 - F(1, 2**41)

# --- smaller Gaussian width: gamma = 2 d alpha^2 <= 8 sqrt(24) d^(3/2) b^(1/2) < 40 b^(25/32)
assert 64*24 < 40**2                       # 8 sqrt(24) < 40
assert 160**32 <= (2**34)**7               # b >= 2^34  ==>  40 b^(25/32) <= b/4  (i.e. b^(7/32) >= 160)

# --- the barrier: g2 < eps * T^2/(1-T) with eps < 1/5 and T < T*(46) ---------------------
Ts = Tstar(b)
print("barrier: (1/5) T*^2/(1-T*) = 2^%.3f" % log2(Ts*Ts/5/(1 - Ts)))
print("all checks passed")
\end{lstlisting}

\begin{thebibliography}{9}
\bibitem{OAI}
OpenAI. \emph{Integer multiplication below $n\log n$}. OpenAI Math
Release preprint OAI:Integer-multiplication-below-n-log-n-September-23-2026,
September 23, 2026.
\url{https://github.com/openai/math}.

\bibitem{HvdH}
D.~Harvey and J.~van der Hoeven. Integer multiplication in time
$O(n\log n)$. \emph{Annals of Mathematics} 193(2):563--617, 2021.
Result numbering as in the 45-page author-hosted manuscript, following
\OAI.

\bibitem{Lean}
L.~de~Moura and S.~Ullrich. The Lean~4 theorem prover and programming
language. In \emph{Automated Deduction -- CADE 28}, LNCS 12699,
pp.~625--635, Springer, 2021.
\end{thebibliography}

\end{document}
